% !TEX root = ../main.tex
% !TEX program = pdflatex
\subsection{Методика измерений}

Все эксперименты реализованы в одной программе \code{src/cache.c} (см. приложение); режим работы задаётся аргументом командной строки. Сборка и запуск всех экспериментов выполняются скриптом \code{run.sh}, графики строятся скриптом \code{plot.py}:
\begin{cmd}
clang -O2 -Wall -o build/cache src/cache.c
./build/cache size 64            > results/size.csv
./build/cache line 256           > results/line_ram.csv
./build/cache line 8             > results/line_l2.csv
./build/cache assoc 128 128 32   > results/assoc_l1.csv
./build/cache assoc 16384 128 32 > results/assoc_l2.csv
python3 plot.py
\end{cmd}

Основной приём~--- \emph{обход цепочки указателей} (pointer chasing). В массиве каждый используемый элемент хранит индекс следующего элемента, и программа выполняет цикл \code{p = a[p]}. Адрес очередного обращения становится известен только после завершения предыдущего, поэтому процессор не может выполнять обращения параллельно, а аппаратная предвыборка (prefetch) не может угадать адрес, если элементы связаны в случайном порядке. В результате время одной итерации равно задержке доступа к тому уровню памяти, в котором находятся данные. Компилятор с ключом \code{-O2} не может удалить такой цикл, так как от каждого чтения зависит следующее, а итоговое значение записывается в \code{volatile}-переменную.

Перед замером выполняется прогревочный проход, затем замер повторяется 5~раз и берётся минимальное время~--- это отсекает влияние прерываний и других процессов. Время измеряется функцией \code{clock_gettime(CLOCK_MONOTONIC_RAW)}.

Процессор Apple~M2 гетерогенный: кэши производительных (P) и энергоэффективных (E) ядер различаются. Средства macOS не позволяют жёстко привязать поток к ядру, поэтому программа устанавливает потоку класс обслуживания \code{QOS_CLASS_USER_INTERACTIVE}, при котором планировщик размещает его на производительных ядрах. Для сравнения в таблице~\ref{tab:sysctl} приведены параметры кэша, которые сообщает ОС (команда \code{sysctl}); они используются только для проверки результатов.

\begin{table}[H]
\centering
\caption{Параметры кэш-памяти Apple M2 по данным \texttt{sysctl}}\label{tab:sysctl}
\begin{tabular}{lcc}
\toprule
Параметр & P-ядра & E-ядра \\
\midrule
L1 кэш данных (на ядро) & 128 КБ & 64 КБ \\
L1 кэш команд (на ядро) & 192 КБ & 128 КБ \\
L2 кэш (общий на кластер из 4 ядер) & 16 МБ & 4 МБ \\
Размер строки (\code{hw.cachelinesize}) & \multicolumn{2}{c}{128 байт} \\
Размер страницы (\code{hw.pagesize}) & \multicolumn{2}{c}{16 КБ} \\
\bottomrule
\end{tabular}
\end{table}

Степень ассоциативности ОС не сообщает, её определение~--- полностью экспериментальная задача.
